\section{Models}
\label{sec:models}

Figure~\ref{fig:overview} shows how the components fit together. All of them
are walk-forward: the model that predicts season $T$ is fitted on seasons
before $T$ only.

\begin{figure}[t]
\centering
\begin{tikzpicture}[
  font=\small,
  comp/.style={draw=black!55, fill=boxfill, rounded corners=2pt, align=center,
               minimum height=10mm, inner sep=3pt, text width=#1},
  comp/.default=29mm,
  res/.style={draw=black!55, fill=white, rounded corners=2pt, align=center,
              minimum height=10mm, inner sep=3pt, text width=25mm},
  arr/.style={-{Stealth[length=1.8mm]}, black!55, shorten >=1pt, shorten <=1pt}]
\node[comp=24mm, minimum height=62mm] (raw) at (0,0) {Play-by-play,\\shifts, shots,\\rosters and\\schedules,\\2008--2026};
\node[comp] (xg)   at (4.4, 2.35)  {Shot quality (xG)\\walk-forward GBM};
\node[comp] (rapm) at (4.4, 0.8)   {RAPM\\stint-level ridge};
\node[comp] (pg)   at (4.4, -0.8)  {Player-game layer\\chained GBMs};
\node[comp] (emb)  at (4.4, -2.35) {Event model and\\player embeddings};
\node[comp=31mm] (season) at (9.0, 2.0)  {Season simulator\\20,000 seasons};
\node[comp=31mm] (h)      at (9.0, 0.0)  {\textbf{\neurhlH{}}\\player layer and\\logistic head with Elo};
\node[comp=31mm] (g)      at (9.0, -2.0) {\textbf{\neurhlG{}}\\single-game engine};
\node[res] (o1) at (13.2, 2.0)  {Standings and\\playoff odds};
\node[res] (o2) at (13.2, 0.0)  {Game win\\probabilities};
\node[res] (o3) at (13.2, -2.0) {Box scores and\\player stat lines};
\foreach \n in {xg, rapm, pg, emb} \draw[arr] (raw.east |- \n) -- (\n.west);
\draw[arr] (xg.east) -- (season.west);
\draw[arr] (rapm.east) -- (season.west);
\draw[arr] (emb.east) -- (h.west);
\draw[arr] (xg.east) -- (g.west);
\draw[arr] (rapm.east) -- (g.west);
\draw[arr] (pg.east) -- (g.west);
\draw[arr, dashed] (h.south) -- node[right, font=\scriptsize] {stacked} (g.north);
\draw[arr] (season.east) -- (o1.west);
\draw[arr] (h.east) -- (o2.west);
\draw[arr] (g.east) -- (o3.west);
\draw[arr] (g.east) -- (o2.west);
\end{tikzpicture}

\caption{Components of \neurhl{}, each fitted walk-forward. RAPM and expected
goals feed the season simulator and \neurhlG{}'s priors and rates; the event
model's player embeddings feed \neurhlH{}; the player-game layer gives
\neurhlG{} its baselines and, summed over the schedule, the player season
projections. \neurhlG{}'s final win probability stacks its own output with
\neurhlH{} and Elo.}
\label{fig:overview}
\end{figure}

\subsection{Baselines}
\label{sec:m-elo}

The primary comparator is the house Elo system (``v1''), tuned on 2006--2017
before any \neurhl{} work began. Before each game with ratings $r_h$ and $r_a$
and home advantage $H$, the expected home score is
\begin{equation}
E_h = \bigl(1 + 10^{-(r_h + H - r_a)/400}\bigr)^{-1}.
\end{equation}
After the game both ratings move by $\pm K M (y - E_h)$, where $y$ is the home
result and $M = \ln(|\Delta g| + 1)\cdot 2.2 / (2.2 + 0.001\,d_w)$ is a
margin-of-victory multiplier that shrinks with the winner's rating advantage
$d_w$. Between seasons each rating regresses toward the mean,
$r \leftarrow \mu + \phi (r - \mu)$. The frozen parameters are $K = 6$,
$H = 35$ and $\phi = 0.7$. Game probabilities come from walk-forward logistic
models of the rating difference: the probability of reaching overtime, the
regulation winner, and the overtime or shootout winner. We also report a
constant home-win rate, and ``Tier-0'', the stronger of a standardised
logistic regression and an early-stopped gradient-boosted model
\citep{chen2016xgboost} on 23
engineered game features (Elo difference, rest, travel, goalie form, roster
aggregates, era), which scored 0.67692 on 2014--2017 against Elo's 0.67548.

\subsection{Shot quality}
\label{sec:m-xg}

The expected-goals model is a histogram gradient-boosted classifier
\citep{friedman2001greedy,ke2017lightgbm} (400
iterations, learning rate 0.06, 31 leaves, at least 200 shots per leaf, L2
penalty 1.0, early stopping) on 33 features: distance and angles, rebound and
rush indicators, the previous event's type, distance, angle and elapsed time,
speed from the previous event, strength state and empty nets, score
difference, period, shooter handedness and position, shot type, and
walk-forward era and rink covariates. For vantage $V$ the model is fitted on
seasons before $V-1$. Its calibration map is isotonic, refitted every 25 league
games on the most recent 120,000 out-of-sample shots strictly before the game
being scored, so it follows the recording drift with the data in which the
drift occurs. The model beats a distance-and-angle logistic regression in all
18 vantages (gate X1), and team expected-goal differential predicts future goal
share better than shot-attempt (Corsi) differential across 554 team-seasons,
$r = 0.423$ against $0.395$ (gate X2). It fails its calibration gate X3, a
maximum deviation of 0.005 per decile, at 0.0065 after three declared
correction ladders; the residual is concentrated in the top decile in the
seasons of the recording change.

\subsection{Adjusted plus-minus}
\label{sec:m-rapm}

RAPM is a stint-level ridge regression \citep{hoerl1970ridge} of on-ice shot rates on indicator
columns for the players on the ice, with an intercept, a home term and
unpenalised team-season fixed effects, so that player coefficients measure
performance relative to the player's own team. The penalty, $\lambda = 51{,}200$,
is chosen to maximise how well a prior rating predicts a player's realised
on-ice results after he changes team. RAPM priors are computed as of each
season and feed the season simulator and \neurhlG{}.

\subsection{Event simulator}
\label{sec:m-events}

The event simulator is a Transformer-Hawkes marked temporal point process
\citep{zuo2020transformer} over the unified event and line-change stream, with
heads for the time to the next event, its type and team, location, zone and
actor, and player effects entering as zero-initialised modulations of the
hazard anchored on RAPM, in the spirit of \citet{thomas2013competing}. It has
5.3 million parameters and is trained as a five-seed ensemble. After the two
leaks described in Section~\ref{sec:failures} were removed it reproduces goals,
shots, penalties and faceoffs per game and per period within 10\% of observed
rates, and its score effects match observed ones with a gradient ratio of
0.979. Its score-effect curve initialises the \neurhlG{} hazard table.

\subsection{Player-game layer}
\label{sec:m-pg}

The player-game layer predicts, for every dressed skater, his share of team
ice time, his shots on goal, and the probabilities of scoring at least one goal
and recording at least one assist. It chains four histogram gradient-boosted
models in the order opportunity, volume, conversion: predicted usage feeds the
shot model, and both feed the two binary models. Declared feature blocks cover
the player's own recent form (the baseline's inputs, nested), usage structure
by strength, line rank and power-play unit, opportunity created by absent
teammates, opponent and schedule context, and line and coaching patterns.
Binary heads are calibrated isotonically on the previous season's
out-of-sample predictions, and ice-time shares are renormalised within each
team-game so they sum to one exactly.

\subsection{Season simulator and player seasons}
\label{sec:m-season}

The season simulator integrates team scoring rates, roster RAPM and score
effects over each game of the real schedule and simulates 20,000 seasons,
drawing each team's strength once per simulated season ($\sigma = 0.07$) so that
uncertainty about how good a team is persists through the season. Player
season projections blend a gradient-boosted season model with the
player-game layer summed over the schedule under an availability hazard,
50/50, the combination a declared rule selected.

\subsection{\neurhlH{}}
\label{sec:m-h}

\neurhlH{} has three layers. Layer~1 predicts what each player will do in a game:
his on-ice shot attempts for and against per 60 minutes, close-range attempts
for and against, and his ice time. It is residual on the player's own
exponentially weighted history: a two-hidden-layer network (128 units, GELU,
dropout 0.1) maps a 64-dimensional player embedding from the event model, 13
static features (position, age and age squared, career games, games with the
current team, a recent-move flag, home, rest, team game number, games played
this season, a late-season flag) and six history features to multiplicative
corrections bounded by $\pm 0.7$ in log space, with a zero-initialised output
layer so that training starts at the baseline. It is trained walk-forward on
about 750,000 player-games. Layer~2 aggregates the predictions over the
players actually dressed, weighted by predicted ice time, into team
projections: projected attempt share and close-range share. Layer~3 is a thin
head, a standardised logistic regression on four inputs, the Elo logit, the
projected attempt-share difference, the projected close-range difference and
the rest difference, refitted for each season on all earlier seasons. Its
regularisation constant is selected on training loss, a known weakness that
favours weak regularisation; it was kept verbatim for the confirmation because
changing it would have tested a different model.

\subsection{\neurhlG{}}
\label{sec:m-g}

\neurhlG{} predicts one game from the two dressed lineups and simulates its box
score. Figure~\ref{fig:g-arch} summarises the architecture.

\paragraph{State.}
After every game each skater's state is updated by discounted sums of his
per-game statistics with per-game decays 0.98, 0.95, 0.85 and 0.60. Rates are
ratios of discounted sums, reported with their effective sample sizes, beside
season-to-date and career totals and as-of priors (the RAPM prior and the
event-model embedding fitted on earlier seasons only). The 114 skater features
cover ice time by strength, individual shots, attempts and expected goals,
goals, assists, on-ice expected goals for and against, and penalties. Goalie
state (23 features) includes shrunk save percentage, goals saved above
expected per expected goal, and workload. Team state (57 features) holds
discounted team rates, and context (10 features) holds the Elo logit, rest,
travel, time-zone change and era covariates. A game late in a season therefore
sees every earlier game of that season as well as all earlier seasons.

\paragraph{Network.}
Shared encoders map each skater (114 features plus position, $\to 128 \to 64$),
each team ($57 \to 64 \to 32$), each goalie ($23 \to 32 \to 16$) and the context
($10 \to 32 \to 16$), each a linear layer, GELU, layer normalisation, dropout
0.2 and a linear layer. Heads produce residual corrections with
zero-initialised output layers: 12 per skater (ice time by strength, shot,
attempt and expected-goal propensities, finishing, playmaking, offensive and
defensive on-ice effects, penalties drawn and taken), six per team and one per
goalie. The untrained network therefore reproduces shrunk-history baselines
exactly. The model has 48,810 parameters.

\paragraph{Team layer.}
The team layer is deterministic and conserves ice time. With $o$ the opponent
of side $s$, power-play opportunities have mean
$\lambda^{pp}_s = \mathrm{log5}(\text{drawn}_s, \text{taken}_o)\,e^{\delta_s}$,
where $\mathrm{log5}(a, b) = ab/L$ for league rate $L$ \citep{hammond2015james}, and $\delta_s$ collects
the skaters' penalty terms. Power-play time is $T^{pp}_s = \tau \lambda^{pp}_s$,
short-handed time is the opponent's power-play time, and even-strength time is
$T^{ev} = 60 - T^{pp}_s - T^{pp}_o$ minutes. Player $i$'s ice time at strength
$k$ is
\begin{equation}
\mathrm{TOI}_{ik} = n_k\, T^{k}\, \mathrm{softmax}_{j \in \mathcal{P}(i)}\bigl(b_{jk} + \delta_{jk}\bigr)_i,
\end{equation}
a softmax over the player's position group $\mathcal{P}(i)$ with $n_k$ players on
the ice (three forwards and two defence at even strength, five skaters on the
power play, four short-handed), so team ice time is exact by construction.
Team scoring rates are log-linear:
\begin{equation}
\log r^{ev}_s = \log \mathrm{log5}\bigl(\mathrm{xGF60}_s, \mathrm{xGA60}_o\bigr)
  + \sum_{i \in s} w_i\, \mathrm{off}_i - \sum_{j \in o} w_j\, \mathrm{def}_j + \delta^{tm}_s,
\end{equation}
where $w$ are ice-time weights and the team rates are shrunk toward league
values by their effective sample sizes. Player shots, attempts, expected goals,
goals, assists and on-ice expected goals are softmax allocations of team
totals, so they sum exactly, and team goals are expected goals times a
finishing term less the opposing goalie's effect.

\paragraph{Outcome.}
Win probabilities come from a differentiable integration over the goal
differential, a Markov chain in the spirit of \citet{kaplan2014markov}. Let $\lambda_h, \lambda_a$ be the two regulation scoring rates. The
distribution of the differential $d \in \{-8, \dots, 8\}$ evolves over 30 steps of
two minutes. In each step side $s$ scores with probability
$q_s = 1 - \exp\{-\lambda_s e^{u}\,\Delta t\, e^{\eta(\ell_s, \pi)}\}$, where
$\ell_s \in \{-3, \dots, 3\}$ is the side's own lead (clamped), $\pi$ indexes five
game phases (minutes 0--40, 40--50, 50--56, 56--58, 58--60), and $\eta$ is a
learned $7 \times 5$ hazard table initialised from the event simulator's score
effects \citep[compare][]{thomas2013competing}. Trailing teams push, leaders
sit back, tied teams turn cautious late, and the last minutes capture pulled
goalies. The pace shock
$u = \sqrt{2}\,\sigma x - \sigma^2/2$ is shared by both teams and integrated with
five-node Gauss--Hermite quadrature \citep{abramowitz1964handbook}, which supplies the tie mass by mechanism
rather than through an ad hoc scalar. The probability of a regulation tie is
the mass at $d = 0$, and the overtime or shootout winner is a logistic function
of $\log(\lambda_h/\lambda_a)$. The output is the four-way outcome: home or away
win, in regulation or after it.

\paragraph{Training.}
Snapshot $T$ is trained on seasons 2009 through $T-1$. Phase~1 minimises
Huber losses on ice time by strength and on power-play minutes, Poisson
deviances on skater and team counts, and the four-way outcome negative log
likelihood with weight 0.3 (AdamW \citep{loshchilov2019decoupled}, learning rate $2\times10^{-3}$, weight decay
0.01, one-cycle schedule, batches of 256 games, 10 epochs). Phase~2 raises
the outcome weight to 1 and trains four more epochs at learning rate
$3\times10^{-4}$, with an L2-SP penalty \citep{li2018explicit} of $10^{-3}$
toward the phase-1 weights. The ensemble averages five seeds.

\paragraph{Final probability.}
For each season $T$ a logistic stack \citep{wolpert1992stacked} over the Elo logit, the \neurhlG{} logit and
the \neurhlH{} logit is fitted on out-of-sample predictions from earlier seasons.
The column subset, which always includes Elo, is chosen by leave-one-season-out
cross-validation on the training seasons, with a fallback to Elo and \neurhlG{}
where \neurhlH{} is unavailable.

\paragraph{Stat sheet.}
Box scores are simulated 10,000 times per game and are consistent with the
win probability by construction. Each draw samples the outcome category from
the stacked four-way probabilities, then a regulation score conditional on it,
then team shots (negative binomial around the model mean, $r = 40$, floored at
goals), team expected goals (gamma), power-play opportunities (Poisson),
player shots and goals (multinomial splits by the model's shares), assists
(zero, one or two per goal in the league's typical mix, split by assist
shares), individual expected goals (Dirichlet split) and goalie saves.

\begin{figure}[t]
\centering
\begin{tikzpicture}[
  font=\footnotesize,
  every node/.append style={execute at begin node={\hyphenpenalty=10000\exhyphenpenalty=10000}},
  b/.style={draw=black!55, fill=boxfill, rounded corners=2pt, align=center,
            minimum height=10mm, inner sep=3pt, text width=#1},
  b/.default=25mm,
  o/.style={draw=black!55, fill=white, rounded corners=2pt, align=center,
            minimum height=10mm, inner sep=3pt, text width=19mm},
  arr/.style={-{Stealth[length=1.6mm]}, black!55, shorten >=1pt, shorten <=1pt},
  lab/.style={font=\scriptsize, text=black!65, inner sep=1.5pt}]
\node[b=24mm] (sk) at (0, 1.35)  {Skaters (both sides)\\114 state features};
\node[b=24mm] (gk) at (0, 0)     {Starting goalies:\\23 features};
\node[b=24mm] (tm) at (0, -1.35) {Team state (57)\\and context (10)};
\node[b=24mm, minimum height=37mm] (enc) at (3.15, 0) {Shared encoders\\and residual heads\\[4pt]{\scriptsize zero-initialised, so\\training starts at\\shrunk baselines}};
\node[b=30mm] (team) at (6.6, 0.9) {Team layer:\\exact ice-time budgets,\\log5 rates, on-ice effects};
\node[b=25mm] (alloc) at (9.9, 0.9) {Allocations: player\\outputs sum to\\team totals};
\node[o] (mc) at (12.85, 0.9) {Monte Carlo\\box scores};
\node[b=28mm] (haz) at (6.55, -1.3) {Hazard integration:\\lead $\times$ phase table,\\Gauss--Hermite pace};
\node[b=25mm] (stk) at (9.9, -1.3) {Logistic stack\\with Elo and \neurhlH{}};
\node[o] (wp) at (12.85, -1.3) {Win\\probability};
\foreach \n in {sk, gk, tm} \draw[arr] (\n.east) -- (enc.west |- \n);
\draw[arr] (enc.east |- team) -- (team.west);
\draw[arr] (team) -- (alloc);
\draw[arr] (alloc) -- (mc);
\draw[arr] (team) -- node[lab, right] {scoring rates} (haz);
\draw[arr] (haz) -- (stk);
\draw[arr] (stk) -- (wp);
\draw[arr] (stk.east) ++(0,2.5mm) -- node[lab, pos=0.62, right] {outcome draw} (mc.south);
\end{tikzpicture}

\caption{\neurhlG{}. Encoders and residual heads adjust shrunk-history baselines; a
deterministic team layer conserves ice time and makes player outputs sum to
team totals; a hazard integration over the goal differential gives the
four-way outcome, which is stacked with Elo and \neurhlH{}; box scores are
simulated conditional on the stacked outcome.}
\label{fig:g-arch}
\end{figure}
